\section{Scope of Price Evidence}
\label{app:scope_adaptation_submission}

% CO-AUTHOR EDIT: expanded thin price appendix into H.1--H.9 (all numbers macro-bound).
The main text (Section~\ref{sec:price_scope_conclusion}) treats price
evidence as corroboration of economic non-neutrality, not as causal
identification, and promises the full price regressions here. This
appendix honors that promise. It records the price sample, the broad
associations, the overlap-cell reweighting that reverses their sign, the
high-value-segment exception, the direct-CADE null, and the boundary
diagnostics that keep the interpretation as scope rather than a price
effect. Every number is the same object reported in the main text; the
full regression grids (all fixed-effect combinations, all quartile cells,
all trim steps) live in the online supplement. The strategic-adaptation
simulation is reported separately in
Appendix~\ref{app:adaptive_deployment_submission}.

\subsection{H.1\quad Price sample and outcome}
\label{app:scope_sample_submission}

The outcome is the log negotiated unit price at the item--award level.
The estimation sample is the broad item-level panel with
$N=\valBetaBroadN$ awards, with item, year, and (where indicated)
procuring-unit (PBU) fixed effects. The same sample underlies every
specification below; the overlap, quartile, and trim variants reweight or
restrict it but do not change the outcome or the unit of observation.

\subsection{H.2\quad Broad price associations}
\label{app:scope_broad_submission}

In the broad specification, frequent-loser presence is positively
associated with the log negotiated unit price. Table~\ref{tab:app_prices_broad_submission}
reports the item$+$year specification ($\valTabPricesColOne$) and the
$+$PBU specification ($\valTabPricesColTwo$), and splits by modality:
Preg\~ao ($\valTabPricesColThree$) and Convite ($\valTabPricesColFour$).
The implied headline range across estimators is $\valHeadlineRange$. The
larger association in the electronic-auction (Preg\~ao) environment is
inconsistent with a pure quorum-filler reading under the three-bidder
rule; these remain descriptive associations, not causal effects.

\begin{table}[htbp]
\centering
\caption{Broad Price Associations (log negotiated unit price)}
\label{tab:app_prices_broad_submission}
\scriptsize
\begin{adjustbox}{max width=\textwidth,center}
\begin{threeparttable}
\begin{tabular}{lcccc}
\toprule
 & Item$+$year & $+$PBU FE & Preg\~ao & Convite \\
 & (1) & (2) & (3) & (4) \\
\midrule
FL present & $\valTabPricesColOne$ & $\valTabPricesColTwo$ & $\valTabPricesColThree$ & $\valTabPricesColFour$ \\
 & $(\valTabPricesSEOne)$ & $(\valTabPricesSETwo)$ & $(\valTabPricesSEThree)$ & $(\valTabPricesSEFour)$ \\
\midrule
$N$ & $\valTabPricesNOne$ & $\valTabPricesNTwo$ & $\valTabPricesNThree$ & $\valTabPricesNFour$ \\
$R^2$ & $\valTabPricesRsqOne$ & $\valTabPricesRsqTwo$ & $\valTabPricesRsqThree$ & $\valTabPricesRsqFour$ \\
\bottomrule
\end{tabular}
\begin{tablenotes}
\scriptsize
\item \textit{Notes:} Log negotiated unit price on a frequent-loser
presence indicator. Columns (1)--(2) pool modalities with item$+$year
and item$+$year$+$PBU fixed effects; (3)--(4) split by modality.
Standard errors clustered at the item level in parentheses. Implied
range across estimators: $\valHeadlineRange$. Descriptive associations,
not causal effects.
$^{**}p<0.05$, $^{***}p<0.01$.
\end{tablenotes}
\end{threeparttable}
\end{adjustbox}
\end{table}

\subsection{H.3\quad Overlap-cell reweighting (sign reversal)}
\label{app:overlap_scope_submission}

The positive broad coefficient is a weighting artifact, not a robust
treatment sign. Holding the overlap cells fixed and moving from
broad-sample weighting to the within-cell average treatment effect on the
treated (ATT) reverses the sign monotonically:
broad $\valBetaBroad$ ($p=\valBetaBroadP$); overlap-cell unweighted
$\valBetaUnwOverlap$ ($p=\valBetaUnwOverlapP$); overlap-cell ATT
$\valScopeOverlapCoef$ ($p\valScopeOverlapPlt$). Propensity-score
trimming pushes it further to $\valScopePSCoef$. Only
$\valDroppedItemsShare$ of treated items lack a within-cell
counterfactual, so the reversal cannot be a dropped-observation artifact.
Table~\ref{tab:app_overlap_decomp_submission} reports the progression.

\begin{table}[htbp]
\centering
\caption{Sign-Reversal Decomposition (overlap-cell reweighting)}
\label{tab:app_overlap_decomp_submission}
\scriptsize
\begin{adjustbox}{max width=\textwidth,center}
\begin{threeparttable}
\begin{tabular}{lcc}
\toprule
Specification & Coefficient & $p$-value \\
\midrule
Broad sample & $\valBetaBroad$ & $\valBetaBroadP$ \\
Overlap-cell, unweighted & $\valBetaUnwOverlap$ & $\valBetaUnwOverlapP$ \\
Overlap-cell ATT & $\valScopeOverlapCoef$ & $\valScopeOverlapPlt$ \\
Overlap-cell ATT, PS-trimmed & $\valScopePSCoef$ & --- \\
\bottomrule
\end{tabular}
\begin{tablenotes}
\scriptsize
\item \textit{Notes:} Same outcome and overlap cells throughout; rows
differ only in weighting and trimming. The broad-to-ATT move reverses
the sign because it strips the between-cell selection component. Only
$\valDroppedItemsShare$ of treated items lack a within-cell
counterfactual.
\end{tablenotes}
\end{threeparttable}
\end{adjustbox}
\end{table}

\subsection{H.4\quad High-value segment and modality heterogeneity}
\label{app:scope_segment_submission}

The negative ATT reading survives both modalities (Convite
$\valSignBetaConvATT$, Preg\~ao $\valSignBetaPregATT$) and every
tender-value quartile except the highest. The high-value tail (Q4) is the
only systematically positive cell ($\valBetaQfour$, $p=\valBetaQfourP$),
where coordination costs are spread over larger contracts. Trimming the
most heavily weighted overlap cells does not eliminate the negative
aggregate: it moves from $\valSRTrimNone$ (no trim) to $\valSRTrimTen$
(drop top decile of cell weights) to $\valSRTrimFifty$ (drop top half).
Table~\ref{tab:app_segment_submission} reports the quartile and trim
cells.

\begin{table}[htbp]
\centering
\caption{High-Value Segment and Trim Sensitivity (overlap-cell ATT)}
\label{tab:app_segment_submission}
\scriptsize
\begin{adjustbox}{max width=\textwidth,center}
\begin{threeparttable}
\begin{tabular}{lc@{\hskip 1.2cm}lc}
\toprule
Cell / modality & ATT & Trim step & ATT \\
\midrule
Convite & $\valSignBetaConvATT$ & No trim & $\valSRTrimNone$ \\
Preg\~ao & $\valSignBetaPregATT$ & Drop top $1\%$ & $\valSRTrimOne$ \\
Tender-value Q1 & $\valSignBetaQoneATT$ & Drop top $10\%$ & $\valSRTrimTen$ \\
Tender-value Q4 & $\valSignBetaQfourATT$ & Drop top $50\%$ & $\valSRTrimFifty$ \\
\bottomrule
\end{tabular}
\begin{tablenotes}
\scriptsize
\item \textit{Notes:} Left block: overlap-cell ATT by modality and by
tender-value quartile; Q4 ($\valBetaQfour$, $p=\valBetaQfourP$) is the
only positive cell. Right block: ATT after dropping the most heavily
weighted overlap cells by cell-weight rank. The negative aggregate is
not driven by a small set of high-weight cells.
\end{tablenotes}
\end{threeparttable}
\end{adjustbox}
\end{table}

\subsection{H.5\quad Direct-CADE overlap price check}
\label{app:scope_cade_submission}

Restricting the overlap to items involving firms named as direct
defendants in CADE cases yields a null price estimate
($\valBetaCADEDirect$, $p=\valBetaCADEDirectP$). The non-direct overlap
estimate is $\valSignBetaCADENonDirectATT$. Price regressions therefore
do not recover the legal object of the adjudicated cases: there is no
price base from which a damages figure could be read off the direct
defendants.

\subsection{H.6\quad Why these are not damages or overcharge}
\label{app:scope_notdamages_submission}

Three facts keep the price evidence on the scope side of the line. First,
the broad positive imprint is selection into structurally higher-priced
cells: among non-treated items, mean log price rises from
$\valSelTestQOne$ (Q1) to $\valSelTestQFive$ (Q5) across cell FL-share
quintiles, a $\Delta=\valSelTestDelta$ log-point gap, with a partialled-out
coefficient of $\valSelTestCoefFullFE$ (SE $\valSelTestSEFullFE$,
$N=\valSelTestN$) once all five cell dimensions are absorbed. The
within-cell object (the ATT) removes exactly this component and is
negative. Second, the direct-CADE null (H.5) means there is no price base
tied to the adjudicated conduct. Third, the design holds the procurement
environment fixed rather than against a counterfactual price, so no
overcharge can be computed. These are scope facts about where flagged
environments sit in the price distribution, not damages.

\subsection{H.7\quad Within-cell mechanism (bidder-count boundary)}
\label{app:scope_mechanism_submission}

The within-cell price compression loads on the count of genuine bidders,
not on the frequent-loser count. The within-cell coefficient is
$\valMechWinnerVsRef$ (SE $\valMechWinnerVsRefSE$) and falls to
$\valMechWinnerVsRefControlled$ (SE $\valMechWinnerVsRefControlledSE$, not
significant) once $\log(\text{n\_firms})$ is added. The price movement
therefore operates through the bidder-count channel; the mechanism is
\emph{not} identified by these data, and we do not read a behavioral
mechanism off the price record.

\subsection{H.8\quad Strategic-adaptation diagnostics}
\label{app:scope_adaptation_diag_submission}

We do not re-run the bunching test here. The density check at the
operational threshold ($T=14$) appears in the profile appendix
(Appendix~\ref{app:profile_threshold_submission}): the local density ratio
is $1.06$, indistinguishable from no bunching, because the threshold was
computed ex post and never disclosed, so firms had no target to game
in-sample. The full strategic-adaptation simulation is in
Appendix~\ref{app:adaptive_deployment_submission}. We cross-reference
rather than duplicate.

\subsection{H.9\quad Limitations}
\label{app:scope_limits_submission}

The price evidence is descriptive, environment-specific, and not causal.
It establishes economic non-neutrality of flagged environments,
concentrated in the high-value tail, but it is not a damages estimate, an
overcharge, a markup, or a treatment effect on price. The coefficients are
BEC-specific and should not be imported mechanically. Another agency
should re-validate the ranking against its own enforcement anchors and
price benchmarks rather than reusing these estimates or the BEC cutoff.
